\cotefile{157-2}
\scanpage{83}\mspage{27}
% L1
C'est \unsure{assez} étrange -- on va voir si c'est bien vrai !

% L2
III \uline{Construction de} $\Phi_0,\Psi_0$ :
\diagram{\begin{tikzcd}[ampersand replacement=\&,column sep=small,row sep=small]
\& \underline{\mathrm{Fibsc}}_0(S)\\
\mathrm{Cat}_{/S}\arrow[ur]\arrow[dr] \&\\
\& \underline{\mathrm{Cofibsc}}_0(S)
\end{tikzcd}}
% L3-4
On doit avoir
\[
\Phi_0(X)=\Phi r C(X)\qquad \Psi_0(X)=\Psi r C(X),
\]
% L5-8
où \deleted{\unsure{r}} \deleted{\unsure{r}} $r:\underline{\mathrm{Parf}}\,S\longrightarrow\mathrm{Cat}_{/S}$ est l'inclusion canonique. En tant que foncteurs à valeurs dans $\underline{\mathrm{Fibsc}}$ resp. $\underline{\mathrm{Cofibsc}}$, cela définit déjà $\Phi_0,\Psi_0$ par
% L9
\[
\Phi_0=\Phi r C,\quad \Psi_0=\Psi r C,
\]
% L10-14
mais il faut voir que $\Phi_0,\Psi_0$ prennent leurs valeurs \deleted{\unsure{à}}dans les sous-catégories pleines à indice $0$. \deleted{\unsure{Parf} \unsure{pour}} Cela équivaut \deleted{\unsure{donc}} au lemme suivant \deleted{(que je formule pour $\Psi_0$, pour fixer les idées)}\unsure{,|;} appliqué à $Y=rC(X)$ :

% L15-17
\uline{Lemme} Soit $Y$ parfait sur $S$ (ou seulement un Hot-fibré sur $S$). Alors $\Phi(\deleted{\unsure{X}}Y)$ est dans $\underline{\mathrm{Fibsc}}_0\,S$, $\Psi(Y)$ dans $\underline{\mathrm{Cofibsc}}_0\,S$.

% L18-19
Il faut prouver \inserted{(pour $\Psi(Y)$)} que si \deleted{\unsure{g}} $s\longrightarrow t$ est une flèche dans $S$, alors
% L20
\[
\Psi(Y)_s\unsure{\mathord{-}|\longrightarrow}\Psi(Y)_t
\]
% L21-23
est une \unsure{W|w}-équivalence. Or cette flèche s'identifie à $\deleted{\unsure{X}}Y_{/s}\longrightarrow Y_{\deleted{\unsure{\vert}}/t}$ et on sait bien que c'est \deleted{un \unsure{Ob|OK}, par définition} \inserted{une Hot-équi}%
% archive sheet number: 82
